%=======================================================================
%  Informe de la Entrega 1 - Proyecto de Integracion Continua (CI/CD)
%  Compilar con:  tectonic informe_entrega1.tex
%  o en Overleaf:  https://www.overleaf.com  (subir esta carpeta)
%=======================================================================
\documentclass[12pt,a4paper]{article}

%-----------------------------------------------------------------------
%  PAQUETES
%-----------------------------------------------------------------------
\usepackage[utf8]{inputenc}
\usepackage[T1]{fontenc}
\usepackage[spanish,es-nodecimaldot]{babel}
% Desactiva el atajo de comillas de babel: el caracter " se usa literalmente
% en nombres de host (por ejemplo "db") y en fragmentos JSON.
\shorthandoff{"}
\usepackage{lmodern}
\usepackage{microtype}
\usepackage{geometry}
\geometry{top=2.5cm,bottom=2.5cm,left=3cm,right=2.5cm}

\usepackage{graphicx}
\usepackage{float}
\usepackage{caption}
\usepackage{booktabs}
\usepackage{array}
\usepackage{longtable}
\usepackage{tabularx}
\usepackage{multirow}
\usepackage{enumitem}
\usepackage{amsmath}
\usepackage{xcolor}
\usepackage{titlesec}
\usepackage{fancyhdr}
\usepackage{listings}
\usepackage{tikz}
\usetikzlibrary{positioning,arrows.meta,shapes.geometric,fit,backgrounds,calc}
\usepackage[hidelinks,breaklinks]{hyperref}

\usetikzlibrary{babel}

%-----------------------------------------------------------------------
%  DATOS EDITABLES  (cambie SOLO esta seccion para personalizar el informe)
%-----------------------------------------------------------------------
\newcommand{\titulo}{Contenedores comunicados con Docker y pipeline de GitHub Actions}
\newcommand{\titulocorto}{Entrega 1 \textbullet\ Docker Compose + GitHub Actions}
\newcommand{\subtitulo}{Entrega 1 -- Aplicaci\'on Flask con PostgreSQL orquestada por Docker Compose}
\newcommand{\materia}{\'Enfasis Profesional I (Integraci\'on Continua)}
\newcommand{\materiaheader}{\'Enfasis Profesional I}
\newcommand{\institucion}{Instituci\'on Educativa \textcolor{gray}{[completar]}}
\newcommand{\docente}{Edward Reyes Corredor}
\newcommand{\repositorio}{\url{https://github.com/SavitarADSO/-nfasis-Profesional-I-Integraci-n-Continua-}}
\newcommand{\grupo}{Grupo \textcolor{gray}{[completar]}}
\newcommand{\fechatexto}{10 de noviembre de 2026}

% Integrantes: se edita SOLO esta linea.
\newcommand{\integrantes}{%
  Miguel \textcolor{gray}{[Apellido]} & \textcolor{gray}{[Matr\'icula]} &
  \textcolor{gray}{[Correo institucional]} \\[2pt]
  Nombre \textcolor{gray}{[Apellido]} & \textcolor{gray}{[Matr\'icula]} &
  \textcolor{gray}{[Correo institucional]} \\[2pt]
  Nombre \textcolor{gray}{[Apellido]} & \textcolor{gray}{[Matr\'icula]} &
  \textcolor{gray}{[Correo institucional]} \\[2pt]
  Nombre \textcolor{gray}{[Apellido]} & \textcolor{gray}{[Matr\'icula]} &
  \textcolor{gray}{[Correo institucional]}
}

%-----------------------------------------------------------------------
%  ESTILOS
%-----------------------------------------------------------------------
\definecolor{azulCJL}{HTML}{1F4E79}
\definecolor{azulClaro}{HTML}{DEEAF6}
\definecolor{verde}{HTML}{2E7D32}
\definecolor{rojo}{HTML}{B00020}
\definecolor{gris}{HTML}{595959}

\titleformat{\section}{\large\bfseries\color{azulCJL}}{\thesection.}{0.6em}{}
\titleformat{\subsection}{\normalsize\bfseries\color{azulCJL}}{\thesubsection}{0.6em}{}
\titlespacing*{\section}{0pt}{1.2em}{0.5em}

\captionsetup{font=small,labelfont=bf,labelsep=period}
\captionsetup[table]{position=top}

\pagestyle{fancy}
\fancyhf{}
\setlength{\headheight}{14.5pt}
\renewcommand{\headrulewidth}{0.4pt}
\fancyhead[L]{\footnotesize\color{gris}\titulocorto}
\fancyhead[R]{\footnotesize\color{gris}\materiaheader}
\fancyfoot[C]{\small\thepage}

\setlength{\parindent}{0.6cm}
\setlength{\parskip}{0.35em}
\sloppy
\setlist[itemize]{leftmargin=1.5em,itemsep=0.2em,topsep=0.3em}
\setlist[enumerate]{leftmargin=1.8em,itemsep=0.2em,topsep=0.3em}

%-----------------------------------------------------------------------
%  LISTINGS
%-----------------------------------------------------------------------
\lstdefinestyle{estilopy}{
  backgroundcolor=\color{azulClaro},
  basicstyle=\ttfamily\footnotesize,
  keywordstyle=\color{azulCJL}\bfseries,
  commentstyle=\color{verde}\itshape,
  stringstyle=\color{rojo},
  numbers=left,
  numberstyle=\tiny\color{gris},
  stepnumber=1,
  numbersep=8pt,
  frame=single,
  rulecolor=\color{gris},
  breaklines=true,
  breakatwhitespace=false,
  tabsize=2,
  showstringspaces=false,
  columns=fullflexible,
  keepspaces=true,
  upquote=true,
  xleftmargin=1.2em
}
\lstdefinestyle{estiloshell}{
  backgroundcolor=\color{black!3},
  basicstyle=\ttfamily\footnotesize,
  keywordstyle=\color{azulCJL}\bfseries,
  commentstyle=\color{verde}\itshape,
  stringstyle=\color{rojo},
  numbers=left,
  numberstyle=\tiny\color{gris},
  numbersep=8pt,
  frame=single,
  rulecolor=\color{gris},
  breaklines=true,
  columns=fullflexible,
  keepspaces=true,
  xleftmargin=1.2em
}
\lstset{style=estilopy}

%-----------------------------------------------------------------------
%  MACRO DE CAPTURAS
%  Si el archivo PNG no existe, se dibuja un recuadro de "captura pendiente"
%  para que el documento compile de todos modos.
%-----------------------------------------------------------------------
\newcommand{\anchofiguraval}{0.92\textwidth}
\newcommand{\anchofigura}[1]{\renewcommand{\anchofiguraval}{#1}}
\newcommand{\captura}[3][]{%
  \begin{figure}[H]
    \centering
    \IfFileExists{#3}{%
      \includegraphics[width=\anchofiguraval]{#3}%
    }{%
      \fbox{%
        \parbox[c][5.0cm][c]{0.88\textwidth}{%
          \centering\color{gris}%
          \textbf{\large EVIDENCIA PENDIENTE}\\[0.6em]
          \small Captura de pantalla a insertar en este punto.\\[0.4em]
          \texttt{\small #3}\\[0.4em]
          \small Reemplace el archivo \texttt{\small #3} en \texttt{figuras/} y recompila.
        }%
      }%
    }%
    \caption{#2}%
    \ifx\relax#1\relax\else\label{#1}\fi
  \end{figure}%
}

\begin{document}

%=======================================================================
%  PORTADA
%=======================================================================
\begin{titlepage}
\centering

\vspace*{1.2cm}
{\color{azulCJL}\rule{\textwidth}{1.6pt}}\\[0.5cm]
{\LARGE\bfseries\color{azulCJL} \titulo\par}
\vspace{0.5cm}
{\large \subtitulo\par}
\vspace{0.5cm}
{\color{azulCJL}\rule{\textwidth}{1.6pt}}

\vspace{1.4cm}
{\color{azulCJL}\large\bfseries Integrantes del equipo\par}
\vspace{0.4cm}
\begin{minipage}{0.86\textwidth}
\small
\renewcommand{\arraystretch}{1.3}
\begin{tabularx}{\textwidth}{>{\bfseries}X >{\bfseries}X}
\textbf{Nombre} & \textbf{Matr\'icula / Correo} \\
\midrule
Miguel Angel Aguilar Luna & SavitarADSO -- \texttt{pee7491@gmail.com} \\
Nombre \textcolor{gray}{[Apellido]} & \textcolor{gray}{[Matr\'icula] -- [correo]} \\
Nombre \textcolor{gray}{[Apellido]} & \textcolor{gray}{[Matr\'icula] -- [correo]} \\
Nombre \textcolor{gray}{[Apellido]} & \textcolor{gray}{[Matr\'icula] -- [correo]} \\
\end{tabularx}
\end{minipage}

\vfill

\begin{minipage}{0.86\textwidth}
\centering
\renewcommand{\arraystretch}{1.3}
\begin{tabularx}{\textwidth}{>{\bfseries}l X}
Materia      & \materia \\
Instituci\'on & \institucion \\
Docente      & \docente \\
Grupo        & \grupo \\
Fecha        & \fechatexto \\
Repositorio  & \repositorio \\
\end{tabularx}
\end{minipage}

\vspace{1.2cm}
{\small\color{gris} Documento generado con LaTeX \quad$\vert$\quad
Proyecto de software basado en herramientas de integraci\'on continua \quad$\vert$\quad
Docker + Docker Compose + GitHub Actions}

\vspace{0.4cm}
\end{titlepage}

%=======================================================================
\tableofcontents
\newpage
\listoffigures
\newpage

%=======================================================================
\section{Resumen}

Este informe documenta la \textbf{Entrega 1} del proyecto de integraci\'on continua de la
materia \textit{Integraci\'on Continua y Entrega Continua (CI/CD)}. La entrega consiste en
construir una aplicaci\'on web sencilla que se ejecuta dentro de un contenedor Docker y que
almacena y consulta informaci\'on en una base de datos PostgreSQL alojada en un segundo
contenedor. Ambos contenedores se definen, configuran y levantan de manera reproducible
mediante un archivo \texttt{docker-compose.yml}, y la automatizaci\'on de la construcci\'on
y las pruebas se realiza con un pipeline de \emph{GitHub Actions}.

El punto t\'ecnico central de la entrega es la \textbf{comunicaci\'on entre contenedores por
nombre de servicio}. En lugar de direccionar la aplicaci\'on hacia \texttt{localhost} o
hacia una direcci\'on IP fija, el c\'odigo resuelve el host \texttt{db}, que es el nombre
de servicio declarado en Docker Compose. Esta decisi\'on hace que la soluci\'on sea
independiente de la topolog\'a\'a de la red, portable entre distintos host y coherente con
el modelo declarativo de Compose. Adem\'as, se implement\'o una gesti\'on responsable de
secretos mediante un archivo \texttt{.env} excluido del control de versiones y una
plantilla versionada \texttt{.env.example}.

Se concludes que la automatizaci\'on del entorno de ejecuci\'on es tan importante como la
automatizaci\'on del c\'odigo: el pipeline construido en esta entrega construye las
imágenes, levanta el stack completo, verifica de extremo a extremo que la aplicaci\'on
guarda y lee datos de la base de datos y desmonta el entorno, garantizando que cada
commit entregue sea ejecutable.

\textbf{Palabras clave:} Docker, Docker Compose, contenedores, PostgreSQL, Flask,
integraci\'on continua, GitHub Actions, redes bridge, gesti\'on de secretos.

%=======================================================================
\section{Introducci\'on}

\subsection{Contexto}

El desarrollo de software moderno exige entregas frecuentes y reproducibles. Un equipo que
depende de que cada miembro instale manualmente las dependencias, configure la base de
datos y pruebe a mano no puede sostener un ritmo de entrega semanal. Docker responde a
esta problem\'atica empaquetando la aplicaci\'on y todas sus dependencias en una imagen
aislada y ligera, mientras que Docker Compose describe en un solo archivo declarativo c\'omo
varios contenedores deben relacionarse entre s\'i.

\subsection{Objetivo general}

Implementar un sistema de dos contenedores comunicados --una API Flask y una base de
datos PostgreSQL-- orquestados con Docker Compose, con la comunicaci\'on resuelta por
nombre de servicio, y automatizar su construcci\'on y verificaci\'on mediante un pipeline
de integraci\'on continua en GitHub Actions.

\subsection{Objetivos espec\'ificos}

\begin{enumerate}
  \item Crear un repositorio de Git con una estructura clara y con \emph{commits} frecuentes
        que documenten cada avance.
  \item Construir la imagen de la aplicaci\'on mediante un \texttt{Dockerfile} que use una
        imagen base con versi\'on fija.
  \item Definir los dos servicios, la red propia, el volumen persistente y la condici\'on de
        dependencia en un \texttt{docker-compose.yml}.
  \item Implementar los endpoints de escritura y lectura sobre PostgreSQL.
  \item Demostrar experimentalmente que la aplicaci\'on se conecta a la base de datos
        usando el nombre de servicio \texttt{db} y no \texttt{localhost}.
  \item Gestionar los secretos con \texttt{.env} (excluido) y \texttt{.env.example}
        (versionado).
  \item Configurar GitHub Actions para construir, probar y desmontar el stack en cada
        \emph{push} o \emph{pull request}.
  \item Entregar el presente informe en formato PDF con evidencia gr\'afica.
\end{enumerate}

\subsection{Alcance y estructura de las entregas}

El proyecto se desarrolla en tres entregas sobre el mismo repositorio:

\begin{center}
\begin{tabularx}{\textwidth}{c c X X}
\toprule
\textbf{Entrega} & \textbf{Semana} & \textbf{Actividad} & \textbf{Incidencia en el repositorio} \\
\midrule
1 & 3 & Dos contenedores comunicados entre s\'i & App Flask, PostgreSQL, \texttt{docker-compose.yml}, workflow de GitHub Actions y este informe. \\
2 & 5 & Jenkins como servidor de integraci\'on continua & Pipeline de Jenkins y documento de requisitos para su implementaci\'on. \\
3 y sustentaci\'on & 7--8 & Todo integrado (contenedores, Jenkins, Travis CI y Codeship) & Documento final con historial de cambios, soluci\'on de problemas, responsabilidades y opiniones. \\
\bottomrule
\end{tabularx}
\end{center}

Este documento cubre \textbf{exclusivamente la Entrega 1}.

%=======================================================================
\section{Marco te\'orico}

\subsection{Contenedores y flujo de construcci\'on de una imagen}

Un contenedor es una instancia ejecutable de una imagen, aislada del sistema anfitri\'on
mediante tecnolog\'ias denamespaces y cgroups del n\'ucleo Linux (Docker Desktop lo
implementa sobre una m\'aquina virtual ligera en Windows y macOS). La ventaja para el
proyecto no es la ``eficiencia'' del contenedor en s\'i, sino la \textbf{reproducibilidad}: la
imagen encapsula el sistema operativo m\'inimo, el int\'erprete y las dependencias exactas
del proyecto.

La construcci\'on de una imagen se define en un \texttt{Dockerfile}, un archivo declarativo
de instrucciones. Cada instrucci\'on genera una capa (layer) que se cachea: si una capa no
cambia, Docker reutiliza la que ya tiene. De ah\'i la pr\'actica de copiar primero el
archivo de dependencias y ejecutar \texttt{pip install} antes de copiar el c\'odigo fuente,
de modo que una modificaci\'on en \texttt{app.py} no invalide la capa de dependencias.

\subsection{Orquestaci\'on con Docker Compose}

Docker Compose es una herramienta que declara un conjunto de servicios en un archivo YAML y
los ejecuta de forma conjunta. En lugar de una secuencia de comandos \texttt{docker run}
complicada de reproducir, se declara el estado deseado y Compose lo alcanza. Entre los
elementos usados en esta entrega se encuentran:

\begin{itemize}
  \item \textbf{Servicios}: cada contenedor del sistema. En este proyecto, \texttt{app}
        (construida desde el \texttt{Dockerfile}) y \texttt{db} (imagen oficial
        \texttt{postgres:16.4}).
  \item \textbf{Nombres de servicio}: alias DNS internos de la red de Compose. Es el
        mecanismo central del paso 5.
  \item \textbf{Dependencias condicionales}: \texttt{depends\_on} con
        \texttt{condition: service\_healthy} hace que el contenedor \texttt{app} no arranque
        hasta que el \emph{health check} de PostgreSQL responda correctamente.
  \item \textbf{Volumenes}: \texttt{db\_data} monta el directorio de datos de PostgreSQL
        (\texttt{/var/lib/postgresql/data}), de modo que la informaci\'on sobrevive a
        \texttt{docker compose down}.
  \item \textbf{Variables de entorno}: \texttt{env\_file} inyecta en ambos contenedores el
        contenido de \texttt{.env}, y el bloque \texttt{environment} fija \texttt{DB\_HOST: db}.
\end{itemize}

\subsection{Redes de Docker y resoluci\'on de nombres}

Compose crea una red \emph{bridge} dedicada al proyecto. Todos los contenedores conectados
a ella comparten un espacio de direcciones y, sobre todo, un \textbf{servidor de DNS
embebido} que resuelve autom\'aticamente el nombre de cada servicio. Gracias a esto, el
contenedor \texttt{app} puede "llamar" a \texttt{db:5432} sin conocer la direcci\'on IP
asignada, que adem\'as puede cambiar cada vez que el stack se levanta.

Esta resoluci\'on por nombre es la alternativa recomendada frente a usar \texttt{localhost}
o una IP fija. Con \texttt{localhost}, la aplicaci\'on intentar\'ia conectarse al propio
contenedor \texttt{app} (el bucle local de cada contenedor es el suyo propio) y la
comunicaci\'on fallar\'ia; con una IP fija, el despliegue se volver\'ia frágil ante
cualquier cambio de topolog\'ia, y obligar\'ia a reconstruir la configuraci\'on en cada
host.

\subsection{Integraci\'on continua con GitHub Actions}

La integraci\'on continua (CI) es la pr\'actica de verificar autom\'aticamente cada cambio
en el c\'odigo mediante la ejecuci\'n de pruebas en un entorno limpio y aislado.
GitHub Actions implementa este modelo con \emph{workflows}, archivos YAML ubicados en
\texttt{.github/workflows/} que describen los pasos (jobs) a ejecutar cuando ocurren eventos
como un \emph{push} o un \emph{pull request}. En esta entrega, el workflow construye las
im\'agenes, levanta el stack, verifica la comunicaci\'on entre contenedores con
\texttt{curl} y finalmente destruye el entorno con \texttt{docker compose down -v},
garantizando que cada ejecuci\'on sea independiente de la anterior.

\subsection{Gesti\'on de secretos}

Las variables de entorno son el mecanismo est\'andar para inyectar configuraci\'on sensible
--como la clave de la base de datos-- sin incrustarla en el c\'odigo ni en la imagen
Docker. En este proyecto se combinan dos archivos: \texttt{.env}, que contiene los valores
reales y est\'a excluido mediante \texttt{.gitignore}, y \texttt{.env.example}, que se
versiona y documenta qu\'e variables existen. El pipeline no necesita conocer la clave
real: genera un \texttt{.env} de prueba a partir de la plantilla en cada ejecuci\'on.

%=======================================================================
\section{Metodolog\'ia}

\subsection{Estructura del repositorio}

\begin{center}
\small
\begin{tabularx}{\textwidth}{@{}>{\small\ttfamily\raggedright\arraybackslash}p{4.4cm} >{\small}X >{\small}c@{}}
\toprule
\textbf{Elemento} & \textbf{Funci\'on} & \textbf{Paso de la r\'ubrica} \\
\midrule
\texttt{app/app.py} & API Flask: escritura, lectura y consulta de mensajes & 4 y 5 \\
\texttt{app/requirements.txt} & Dependencias con versi\'on fijada (\textit{pin}) & 2 \\
\texttt{app/Dockerfile} & Imagen de la aplicaci\'on con base \texttt{python:3.12-slim} & 2 \\
\texttt{app/.dockerignore} & Excluye \texttt{.env}, cach\'e y \texttt{.git} del build & 2 y 6 \\
\texttt{docker-compose.yml} & Servicios, red, volumen, \texttt{depends\_on} y \emph{health check} & 3 \\
\texttt{.env.example} & Plantilla de variables de entorno versionada & 6 \\
\texttt{.gitignore} & Excluye \texttt{.env}, cach\'e de Python e IDE & 6 \\
\texttt{.github/workflows/}\newline\texttt{docker.yml} & Pipeline de construcci\'on y pruebas & 7 \\
\texttt{README.md} & Instructivo de uso, comandos y soluci\'on de problemas & 1 y 8 \\
\texttt{docs/} & Informe en LaTeX y capturas de evidencia & 8 \\
\bottomrule
\end{tabularx}
\end{center}

\subsection{Diagrama de arquitectura}

\begin{figure}[H]
\centering
\begin{tikzpicture}[
  font=\small,
  escala/.style={draw=azulCJL,fill=azulClaro,rounded corners=3pt,minimum width=2.9cm,minimum height=1.5cm,align=center,thick},
  cliente/.style={draw=gris,fill=black!5,rounded corners=3pt,minimum width=2.3cm,minimum height=1.3cm,align=center},
  vol/.style={draw=verde,fill=verde!12,rounded corners=3pt,minimum width=2.8cm,minimum height=1.1cm,align=center},
  ci/.style={draw=rojo,fill=rojo!10,rounded corners=3pt,minimum width=3.1cm,minimum height=1.3cm,align=center},
  flecha/.style={-{Latex[length=2.6mm]},thick,color=azulCJL},
  etiqueta/.style={font=\scriptsize,align=center,inner sep=2pt}
]
  % Nodos
  \node[cliente] (cli) {Navegador\\o \texttt{curl}};
  \node[escala,right=2.1cm of cli] (appc) {\textbf{Contenedor 1}\\\texttt{app}\\API Flask\\(\texttt{python:3.12-slim})};
  \node[escala,right=2.4cm of appc] (dbc) {\textbf{Contenedor 2}\\\texttt{db}\\PostgreSQL 16.4};
  \node[vol,below=1.0cm of dbc] (vol) {Volumen \texttt{db\_data}\\$\texttt{/var/lib/postgresql/data}$};
  \node[ci,above=1.6cm of appc] (gh) {\textbf{GitHub Actions}\\\texttt{build} $\rightarrow$ \texttt{up} $\rightarrow$ \texttt{test}};

  % Aristas
  \draw[flecha] (cli) -- node[etiqueta,above=8mm] {TCP 5000\\\texttt{localhost:5000}} (appc);
  \draw[flecha] (appc) -- node[etiqueta,above=8mm] {red \texttt{red\_proyecto}\\host \texttt{"db"}:5432} (dbc);
  \draw[flecha] (dbc) -- (vol);
  \draw[flecha,dashed,color=rojo] (gh) -- node[etiqueta,right=2mm,align=left] {\texttt{docker compose}\\\texttt{build / up / down}} (appc);

  % Etiqueta de red
  \node[draw=azulCJL,dashed,rounded corners=4pt,fit=(appc)(dbc)(vol),inner sep=4mm,opacity=0.25] {};
\end{tikzpicture}
\caption{Arquitectura de la soluci\'on: la aplicaci\'on Flask (contenedor \texttt{app}) accede
a PostgreSQL (contenedor \texttt{db}) a trav\'es de la red propia de Compose usando el
nombre de servicio \texttt{db}; los datos persisten en el volumen \texttt{db\_\allowbreak data}.}
\label{fig:arquitectura}
\end{figure}

\subsection{Flujo de datos de una operaci\'on}

\begin{enumerate}
  \item El cliente env\'ia \texttt{POST /mensajes} con un JSON \texttt{\{"texto":"hola"\}} al
        puerto 5000 publicado del contenedor \texttt{app}.
  \item Flask valida la carga y abre una conexi\'on a PostgreSQL usando
        \texttt{host=DB\_HOST}, cuyo valor es \texttt{db}.
  \item El servidor DNS interno de la red de Compose resuelve \texttt{db} a la direcci\'on
        IP del contenedor de base de datos.
  \item La instrucci\'on \texttt{INSERT} se ejecuta y PostgreSQL escribe la fila en el
        volumen \texttt{db\_data}.
  \item Un \texttt{GET /mensajes} recorre el camino inverso: la app consulta la tabla y
        devuelve el JSON con los registros.
\end{enumerate}

%=======================================================================
\section{Implementaci\'on: desarrollo de los ocho pasos}

\subsection{Paso 1. Repositorio de Git con commits frecuentes}

Se cre\'o el repositorio \texttt{proyecto-ci} y se resolvi\'o evitar un \emph{commit} \'unico
que contenga todo el proyecto. Cada componente del enunciado se a\'adi\'o en un commit
independiente con un mensaje descriptivo, lo que permite auditar la evoluci\'on del
proyecto:

\begin{center}
\small
\begin{tabularx}{\textwidth}{@{}>{\small\ttfamily\raggedright\arraybackslash}p{5.4cm} >{\small}X@{}}
\toprule
\textbf{Commit} & \textbf{Contenido} \\
\midrule
\texttt{Agrega aplicaci\'on Flask y dependencias} & \texttt{app/app.py}, \texttt{app/requirements.txt} \\
\texttt{Agrega Dockerfile y .dockerignore} & Imagen de la app con base de versi\'on fija \\
\texttt{Agrega docker-compose con red y volumen} & \texttt{docker-compose.yml} \\
\texttt{Agrega gesti\'on de secretos} & \texttt{.env.example} y \texttt{.gitignore} \\
\texttt{Agrega pipeline de GitHub Actions} & \texttt{.github/workflows/docker.yml} \\
\texttt{Agrega README con instrucciones del proyecto} & Documentaci\'on del proyecto \\
\texttt{Agrega informe LaTeX y capturas de evidencia} & \texttt{docs/informe\_entrega1.tex}, \texttt{docs/figuras/} \\
\texttt{Actualiza los datos de la portada del informe} & Datos de la portada y resultados de las pruebas \\
\texttt{Actualiza la evidencia del historial de commits} & Captura del \texttt{git log} y figura del pipeline \\
\texttt{Ajustes finales del informe} & Versionado del PDF y correcciones de maquetaci\'on \\
\bottomrule
\end{tabularx}
\end{center}

\begin{lstlisting}[caption=Historial de commits del repositorio,style=estiloshell]
$ git log --oneline
fbbd892 Agrega README con instrucciones del proyecto
31cda33 Agrega pipeline de GitHub Actions
58379cb Agrega gestion de secretos (.env.example y .gitignore)
14cfaff Agrega docker-compose con red propia y volumen
4ed944a Agrega Dockerfile y .dockerignore
7b86dc1 Agrega aplicacion Flask y dependencias
\end{lstlisting}

Adem\'as de los seis commits anteriores, se realizaron tres commits finales con el
informe en LaTeX, las capturas de evidencia y los datos de la portada. De este modo el
historial completo del repositorio queda documentado y cada commit refleja un avance
concreto del proyecto.

\anchofigura{0.75\textwidth}
\captura[fig:gitlog]{Historial de commits del repositorio en GitHub, donde se observa la
secuencia de avances independientes.}{figuras/01-git-log.png}
\anchofigura{0.92\textwidth}

\subsection{Paso 2. Dockerfile con imagen base de versi\'on fija}

El \texttt{Dockerfile} se apoya en \texttt{python:3.12-slim}, una etiqueta con versi\'on
expl\'icita. No se usa \texttt{python:latest} porque una etiqueta m\'vil rompe la
reproducibilidad: la misma instrucci\'on de construcci\'on podr\'ia producir im\'agenes
distintas en meses diferentes, y una actualizaci\'on autom\'atica del sistema base podr\'ia
introducir regresiones no detectadas.

\begin{lstlisting}[caption=Contenido de app/Dockerfile]
FROM python:3.12-slim

ENV PYTHONDONTWRITEBYTECODE=1
ENV PYTHONUNBUFFERED=1

WORKDIR /app

# Primero las dependencias (aprovecha la cache de capas)
COPY requirements.txt .
RUN pip install --no-cache-dir -r requirements.txt

# Despues el codigo
COPY . .

EXPOSE 5000

CMD ["python", "app.py"]
\end{lstlisting}

El orden de las instrucciones \texttt{COPY} es deliberado: al copiar primero
\texttt{requirements.txt} y ejecutar \texttt{pip install}, la capa costosa de instalaci\'on
se cachea. Un cambio en el c\'odigo de la aplicaci\'on solo invalida la \'ultima capa.

\subsection{Paso 3. Orquestaci\'on con Docker Compose}

El archivo \texttt{docker-compose.yml} declara los dos servicios, la red propia, el volumen
persistente y la condici\'on de dependencia con \emph{health check}. Esta \'ultima es
importante: PostgreSQL acepta conexiones unos segundos despu\'es de arrancar el
contenedor, por lo que \texttt{depends\_on} simple no garantiza que la app encuentre la base
de datos lista. El \emph{health check} basado en \texttt{pg\_isready} resuelve el problema.

\begin{lstlisting}[caption=Extracto de docker-compose.yml,language={}]
services:
  app:
    build: ./app
    container_name: ci_app
    ports:
      - "5000:5000"
    env_file: .env
    environment:
      DB_HOST: db          # nombre del servicio, NO localhost
      DB_PORT: 5432
    depends_on:
      db:
        condition: service_healthy
    networks:
      - red_proyecto

  db:
    image: postgres:16.4
    container_name: ci_db
    env_file: .env
    volumes:
      - db_data:/var/lib/postgresql/data
    healthcheck:
      test: ["CMD-SHELL", "pg_isready -U $${POSTGRES_USER} -d $${POSTGRES_DB}"]
      interval: 5s
      timeout: 3s
      retries: 10
    networks:
      - red_proyecto

networks:
  red_proyecto:
    driver: bridge

volumes:
  db_data:
\end{lstlisting}

\subsection{Paso 4. Aplicaci\'on Flask y base de datos}

La aplicaci\'on expone cuatro endpoints: un verificador de salud que incluye el estado de
la conexi\'on a la base de datos, la creaci\'on de mensajes, el listado completo y la
consulta individual por identificador.

\begin{lstlisting}[caption=Fragmento de app/app.py: conexion por nombre de servicio,language=Python]
DB_HOST = os.environ.get("DB_HOST", "db")  # "db" = nombre del servicio, NO localhost

def get_conn():
    """Abre una conexion nueva a PostgreSQL usando el host por nombre de servicio."""
    return psycopg2.connect(
        host=DB_HOST,
        port=DB_PORT,
        dbname=POSTGRES_DB,
        user=POSTGRES_USER,
        password=POSTGRES_PASSWORD,
        connect_timeout=5,
    )

@app.post("/mensajes")
def crear_mensaje():
    datos = request.get_json(silent=True) or {}
    texto = (datos.get("texto") or "").strip()
    if not texto:
        return jsonify(error="El campo 'texto' es obligatorio"), 400
    with conexion() as cur:
        cur.execute("INSERT INTO mensajes (texto) VALUES (%s) RETURNING id, texto", (texto,))
        nuevo_id, texto_guardado = cur.fetchone()
    return jsonify(id=nuevo_id, texto=texto_guardado), 201
\end{lstlisting}

Los par\'ametros de las consultas se env\'ian con marcadores (\texttt{\%s}) de psycopg2 y
nunca por concatenaci\'on de cadenas, lo que evitainyecci\'on SQL.

\subsection{Paso 5. Comunicaci\'on entre contenedores por nombre de servicio}

Este es el paso de mayor peso en la r\'ubrica (35 de 150 puntos). La evidencia se compone de
cuatro pruebas complementarias:

\begin{enumerate}
  \item \textbf{Escritura (app $\rightarrow$ db)}: un \texttt{POST /mensajes} inserta una
        fila y devuelve el identificador generado por PostgreSQL.
  \item \textbf{Lectura (db $\rightarrow$ app)}: un \texttt{GET /mensajes} devuelve la fila
        previamente almacenada, lo que demuestra la persistencia real, no una simulaci\'on.
  \item \textbf{Resoluci\'on de nombres}: dentro del contenedor \texttt{app} se resuelve
        \texttt{db} con \texttt{socket.\allowbreak gethostbyname}, que devuelve la IP interna del
        contenedor de base de datos y confirma que no hay \texttt{localhost} ni IP fija
        escrita a mano.
  \item \textbf{Inspecci\'on desde la base de datos}: un \texttt{SELECT} en \texttt{psql}
        muestra la misma fila, probando que el dato lleg\'o realmente a PostgreSQL.
\end{enumerate}

\begin{lstlisting}[caption=Comandos de verificacion de la comunicacion,style=estiloshell]
# 1. Guardar un dato (app -> db)
curl -X POST http://localhost:5000/mensajes \
     -H "Content-Type: application/json" -d "{\"texto\":\"hola\"}"

# 2. Leerlo (db -> app)
curl http://localhost:5000/mensajes

# 3. Demostrar que la app resuelve "db" por nombre de servicio
docker compose exec app python -c "import socket; print(socket.gethostbyname('db'))"

# 4. Ver el dato directamente en PostgreSQL
docker compose exec db psql -U usuario -d proyecto -c "SELECT * FROM mensajes;"
\end{lstlisting}

En PowerShell, \texttt{curl} es un alias de \texttt{Invoke-WebRequest} y no admite los
mismos argumentos, por lo que se utiliza el equivalente:

\begin{lstlisting}[caption=Equivalente en PowerShell,style=estiloshell]
Invoke-RestMethod -Method Post -Uri http://localhost:5000/mensajes `
  -ContentType "application/json" -Body '{"texto":"hola"}'
\end{lstlisting}

\anchofigura{0.45\textwidth}
\captura[fig:post]{Escritura de un mensaje en la base de datos a trav\'es de la API
(\texttt{POST /mensajes}); la respuesta incluye el \texttt{id} asignado por PostgreSQL.}{figuras/06-curl-post.png}
\anchofigura{0.45\textwidth}
\captura[fig:get]{Lectura de los mensajes almacenados (\texttt{GET /mensajes}); el dato
ingresado en la figura anterior aparece en la respuesta.}{figuras/07-curl-get.png}
\anchofigura{0.30\textwidth}
\captura[fig:dns]{Resoluci\'on del nombre de servicio \texttt{db} desde dentro del
contenedor \texttt{app}; la direcci\'on devuelta (\texttt{172.18.0.2}) es la IP interna
del contenedor de base de datos dentro de \texttt{red\_proyecto}.}{figuras/08-dns-db.png}
\anchofigura{0.60\textwidth}
\captura[fig:psql]{Consulta \texttt{SELECT} ejecutada directamente en PostgreSQL, que
confirma la persistencia del mensaje en el volumen \texttt{db\_data}.}{figuras/09-psql-select.png}
\anchofigura{0.92\textwidth}

\subsection{Paso 6. Gesti\'on de secretos: \texttt{.env} y \texttt{.env.example}}

El archivo \texttt{.env} contiene los valores reales y est\'a listado en
\texttt{.gitignore}, de modo que jam\'as llega al repositorio. El archivo
\texttt{.env.example} s\'i se versiona y cumple tres funciones: documenta qu\'e variables
existen, permite a cualquier compa\~nero reconstruir su entorno local con un solo comando
y sirve como origen del \texttt{.env} de prueba que genera el pipeline.

\begin{lstlisting}[caption=Contenido de .env.example,language={}]
POSTGRES_DB=proyecto
POSTGRES_USER=usuario
POSTGRES_PASSWORD=cambia_esta_clave
\end{lstlisting}

\begin{lstlisting}[caption=Fragmento de .gitignore,language={}]
# Secretos: NUNCA subir el archivo .env real (paso 6).
.env

# Python
__pycache__/
*.py[cod]
\end{lstlisting}

\begin{lstlisting}[caption=Verificacion de que .env esta ignorado por Git,style=estiloshell]
$ git check-ignore -v .env
.gitignore:3:.env        .env

$ git ls-files | grep -c "\.env$"
1                       # solo .env.example esta versionado
\end{lstlisting}

Un detalle pr\'actico: si se cambia la clave de \texttt{.env} despu\'es de haber levantado
la base de datos una vez, PostgreSQL conserva la clave original porque el volumen
\texttt{db\_data} ya tiene inicializado el cl\'uster. La soluci\'on es ejecutar
\texttt{docker compose down -v} para destruir el volumen y volver a levantar.

\subsection{Paso 7. Pipeline de integraci\'on continua en GitHub Actions}

\begin{lstlisting}[caption=Contenido de .github/workflows/docker.yml,language={}]
name: Docker CI

on:
  push:
    branches: [main]
  pull_request:
    branches: [main]

jobs:
  build:
    runs-on: ubuntu-latest
    steps:
      - uses: actions/checkout@v4
      - name: Crear .env de prueba
        run: cp .env.example .env
      - name: Construir imagenes
        run: docker compose build
      - name: Levantar y probar comunicacion
        run: |
          docker compose up -d
          sleep 20
          curl -f http://localhost:5000/
          curl -f -X POST http://localhost:5000/mensajes \
            -H "Content-Type: application/json" -d '{"texto":"hola desde CI"}'
          curl -f http://localhost:5000/mensajes
          docker compose exec app python -c "import socket; print(socket.gethostbyname('db'))"
      - name: Apagar
        if: always()
        run: docker compose down -v
\end{lstlisting}

Dos detalles relevantes del pipeline: el paso de apagado usa \texttt{if: always()}, de modo
que los contenedores se destruyen aunque una prueba falle, evitando fugas de recursos en
el runner; y la generaci\'on del \texttt{.env} de prueba a partir de \texttt{.env.example}
garantiza que la ejecuci\'on no dependa de secretos reales y que el pipeline valide
exactamente la misma ruta de configuraci\'on que usa un desarrollador local.

\captura[fig:actions]{Ejecuci\'on exitosa del pipeline en la pestaña \emph{Actions} del
repositorio: la corrida \texttt{Docker CI} aparece en verde y los seis pasos ---copiar el
c\'odigo, crear el \texttt{.env} de prueba, construir las im\'agenes, levantar y probar
la comunicaci\'on entre contenedores y apagar el entorno--- aparecen con conclusi\'on
\texttt{success}.}{figuras/10-actions-verde.png}

\subsection{Paso 8. Informe final en PDF}

El presente documento se genera con LaTeX. El c\'odigo fuente se encuentra en
\texttt{docs/informe\_entrega1.tex} y las capturas de pantalla en \texttt{docs/figuras/}.
Las referencias a las figuras se resuelven con \texttt{\textbackslash ref}, por lo que si una
captura se agrega o se mueve, las numeraciones se actualizan solas al recompilar. El
documento se compila localmente con un motor LaTeX (\emph{Tectonic} o \texttt{pdflatex}) o
en l\'inea mediante Overleaf, y se entrega en PDF.

Cada figura del informe se inserta con una macro que comprueba si el archivo de imagen
existe. Si a\'un no se ha capturado una pantalla, el documento muestra un recuadro que
indica el nombre exacto del archivo pendiente y compila sin errores, de modo que nunca se
entrega un PDF roto por una captura faltante. La lista completa de las capturas
solicitadas, con su contenido esperado, se documenta en \texttt{docs/figuras/LEEME.md}.

%=======================================================================
\section{Pruebas y resultados}

\subsection{Caso de prueba de extremo a extremo}

\begin{center}
\small
\begin{tabularx}{\textwidth}{@{}c >{\small}X >{\small}X c@{}}
\toprule
\textbf{ID} & \textbf{Precisi\'on} & \textbf{Resultado esperado} & \textbf{Estado} \\
\midrule
CP-01 & El contenedor \texttt{app} arranca y responde en \texttt{/} & \texttt{200} con \texttt{\{"status":"ok"\}} y \texttt{"db\_estado":"ok"} & Aprobado \\
CP-02 & Se inserta un mensaje con \texttt{POST /mensajes} & \texttt{201} con el \texttt{id} generado por PostgreSQL & Aprobado \\
CP-03 & Se listan los mensajes con \texttt{GET /mensajes} & \texttt{200} con el mensaje insertado en JSON & Aprobado \\
CP-04 & Se consulta un mensaje por \texttt{id} & \texttt{200} con el texto; \texttt{404} si no existe & Aprobado \\
CP-05 & Se resuelve \texttt{db} desde el contenedor \texttt{app} & IP interna \texttt{172.18.0.2} del contenedor \texttt{db} & Aprobado \\
CP-06 & El dato es visible en PostgreSQL con \texttt{psql} & Una fila con el texto insertado & Aprobado \\
CP-07 & \texttt{DB\_HOST} dentro del contenedor & \texttt{DB\_HOST=db} (nunca \texttt{localhost}) & Aprobado \\
CP-08 & \texttt{.env} no aparece en el repositorio & \texttt{git ls-files} solo lista \texttt{.env.example} & Aprobado \\
CP-09 & El pipeline de GitHub Actions se ejecuta completo & Todos los pasos en verde & Aprobado \\
\bottomrule
\end{tabularx}
\end{center}

Los casos CP-01 a CP-08 se ejecutaron localmente contra el stack levantado con
\texttt{docker compose up -d}, y su evidencia se muestra en las figuras de esta secci\'on y
de la secci\'on \ref{fig:post}. El caso CP-09 se valid\'o con la primera ejecuci\'on del
workflow en GitHub Actions, que termin\'o con conclusi\'on \texttt{success} y con los seis
pasos en verde (ver figura \ref{fig:actions}).

\subsection{Evidencia gr\'afica del entorno}

\anchofigura{0.80\textwidth}
\captura[fig:build]{Construcci\'on de las im\'agenes con \texttt{docker compose build
--no-cache}; se observa la exportaci\'on de las capas y la creaci\'on de la imagen
\texttt{proyecto-ci-app}.}{figuras/02-docker-build.png}
\anchofigura{0.92\textwidth}

\anchofigura{0.45\textwidth}
\captura[fig:up]{Levantamiento de los dos contenedores con \texttt{docker compose up -d};
el log muestra la creaci\'on del volumen \texttt{db\_data}, de la red
\texttt{red\_proyecto} y el arranque ordenado de \texttt{ci\_db} (healthy) y
\texttt{ci\_app}.}{figuras/03-docker-compose-up.png}
\anchofigura{0.92\textwidth}

\captura[fig:ps]{Salida de \texttt{docker ps} con los contenedores \texttt{ci\_app} y
\texttt{ci\_db} en estado \emph{Up (healthy)}, y sus puertos e identificadores.}{figuras/04-docker-ps.png}

\anchofigura{0.32\textwidth}
\captura[fig:network]{Salida de \texttt{docker network inspect}: los dos contenedores
pertenecen a la misma red \texttt{red\_proyecto} y cada uno tiene asignada una direcci\'on
IP interna, que es la que resuelve el nombre de servicio desde la app.}{figuras/05-network-inspect.png}
\anchofigura{0.92\textwidth}

\anchofigura{0.55\textwidth}
\captura[fig:logs]{Logs del contenedor \texttt{app}: se observa que la base de datos
est\'a disponible desde el primer intento, que la tabla \texttt{mensajes} se crea al
arrancar y que Flask escucha en \texttt{0.0.0.0:5000}, direcci\'on necesaria para que el
contenedor sea alcanzable desde fuera.}{figuras/10-logs.png}
\anchofigura{0.92\textwidth}

%=======================================================================
\section{Soluci\'on de problemas}

\begin{center}
\small
\begin{tabularx}{\textwidth}{@{}>{\small}p{3.2cm} >{\small}X >{\small}X@{}}
\toprule
\textbf{S\'intoma} & \textbf{Causa} & \textbf{Soluci\'on aplicada} \\
\midrule
La app no encuentra la base de datos &
El contenedor \texttt{app} arranc\'o antes de que PostgreSQL aceptara conexiones &
Se a\~nadi\'o un \emph{health check} con \texttt{pg\_isready} y
\texttt{depends\_on: condition: service\_healthy}; adem\'as la app reintenta la conexi\'on
durante el arranque \\
\texttt{could not translate host name "db"} &
El contenedor qued\'o fuera de la red de Compose o \texttt{DB\_HOST} qued\'o como
\texttt{localhost} &
Se verific\'o con \texttt{docker network inspect} y con
\texttt{docker compose exec app env} \\
Error de autenticaci\'on contra PostgreSQL &
La clave de \texttt{.env} cambi\'o despu\'es de inicializarse el volumen &
Se elimin\'o el volumen con \texttt{docker compose down -v} y se reconstruy\'o el stack \\
El pipeline fallaba de forma intermitente &
El sleep fijo no garantizaba que la base estuviera lista &
Se sustituy\'o la espera fija por la condici\'on \texttt{service\_healthy} y se a\~nadi\'o
un \emph{health check} al servicio \texttt{app} \\
\texttt{curl} no funcionaba en PowerShell &
\texttt{curl} es un alias de \texttt{Invoke-WebRequest} en Windows &
Se document\'o el equivalente con \texttt{Invoke-RestMethod} \\
El puerto 5000 aparec\'a ocupado &
Otro proceso local usaba el mismo puerto &
Se identific\'o con \texttt{netstat -ano} y \texttt{findstr :5000}, y se liber\'o el puerto \\
\bottomrule
\end{tabularx}
\end{center}

%=======================================================================
\section{Conclusiones}

La primera conclusi\'on es que el uso del \textbf{nombre de servicio} en lugar de
\texttt{localhost} o de una direcci\'on IP fija no es un detalle menor de configuraci\'on,
sino el mecanismo que hace posible la comunicaci\'on entre contenedores. Dentro de la red
de Compose, cada servicio recibe un alias DNS gestionado por el servidor de nombres
embebido de Docker; ese alias resuelve a la direcci\'on que el contenedor tenga en ese
momento. Gracias a ello, la aplicaci\'on queda completamente ajena a la topolog\'ia de la
red: el mismo c\'odigo funciona en el port\'atil, en un servidor propio y en el runner de
GitHub Actions, sin una sola l\'inea de configuraci\'on diferente. Adem\'as, usar
\texttt{localhost} dentro de un contenedor es un error conceptual grave, porque el bucle
local de cada contenedor apunta a s\'i mismo y no al contenedor vecino.

La segunda conclusi\'on es que \textbf{Docker Compose aporta la declaratividad y la
reproducibilidad} que la ejecuci\'on manual no ofrece. En lugar de memorizar una decena de
comandos \texttt{docker run} con sus par\'ametros de red, de puertos y de variables de
entorno, se declara en un \'unico archivo el estado deseado del sistema y Compose se
encarga de materializarlo. Las \'ulimas consecuencias pr\'acticas son la declaraci\'on
expl\'icita de la red y del volumen, la condici\'on de dependencia verificada mediante
\emph{health checks} --no un simple \texttt{depends\_on} que solo espera el arranque del
proceso-- y la posibilidad de destruir y reconstruir el entorno completo con un solo
comando, lo que resulta invaluable en el pipeline, donde cada ejecuci\'on debe partir de un
estado limpio.

La tercera conclusi\'on tiene que ver con la \textbf{gesti\'on de secretos} y con su
efecto en la calidad del software. El uso del par \texttt{.env} (ignorado por Git) y
\texttt{.env.example} (versionado) separa la configuraci\'on sensible de la
configuraci\'on compartida: la clave de la base de datos nunca queda escrita en el
c\'odigo ni en la imagen de Docker, y a la vez cualquier compa\~nero puede reconstruir su
entorno local con un \'unico comando a partir de la plantilla. Esta separaci\'on es la que
permite, adem\'as, que el pipeline genere su propio \texttt{.env} de prueba y valide la
misma ruta de c\'odigo que se usa en desarrollo.

Finalmente, la integraci\'on continua result\'o ser el multiplicador del trabajo
realizado. El workflow construido no se limita a construir las im\'agenes: levanta el
stack completo, ejecuta la prueba de escritura y lectura con \texttt{curl}, comprueba que
el nombre \texttt{db} se resuelve correctamente y desmonta todo al terminar. En otras
palabras, el pipeline verifica el mismo comportamiento que se documenta en este informe, lo
que convierte las capturas de pantalla en evidencia reproducible y no en unasimple
fotograf\'ia de un momento. El desaf\'io de las entregas siguientes ---incorporar Jenkins,
Travis CI y Codeship sobre esta base--- parte de una fundaci\'on ya verificada, y la
lecci\'on principal es que la automatizaci\'on del entorno de ejecuci\'on es tan
determinante como la automatizaci\'on del c\'odigo mismo.

%=======================================================================
\section*{Referencias}
\addcontentsline{toc}{section}{Referencias}

\begin{list}{}{\setlength{\leftmargin}{1.5cm}\setlength{\itemsep}{0.5em}}
\item Docker Inc. (s.\ f.). \textit{Compose specification}. Recuperado el 10 de noviembre
      de 2026, de \url{https://docs.docker.com/compose/compose-file/}

\item Docker Inc. (s.\ f.). \textit{Docker Compose overview}. Recuperado el 10 de noviembre
      de 2026, de \url{https://docs.docker.com/compose/}

\item Docker Inc. (s.\ f.). \textit{Networking in Compose}. Recuperado el 10 de noviembre
      de 2026, de \url{https://docs.docker.com/compose/how-tos/networking/}

\item Docker Inc. (s.\ f.). \textit{Use volumes}. Recuperado el 10 de noviembre de 2026,
      de \url{https://docs.docker.com/engine/storage/volumes/}

\item Docker Inc. (s.\ f.). \textit{Dockerfile reference}. Recuperado el 10 de noviembre de
      2026, de \url{https://docs.docker.com/reference/dockerfile/}

\item GitHub (s.\ f.). \textit{About GitHub Actions}. Recuperado el 10 de noviembre de 2026,
      de \url{https://docs.github.com/en/actions}

\item GitHub (s.\ f.). \textit{Workflow syntax for GitHub Actions}. Recuperado el 10 de
      noviembre de 2026, de \url{https://docs.github.com/en/actions/reference/workflows-and-actions/workflow-syntax}

\item Python Software Foundation. (2024). \textit{Flask documentation (3.0.x)}.
      Recuperado el 10 de noviembre de 2026, de \url{https://flask.palletsprojects.com/en/3.0.x/}

\item PostgreSQL Global Development Group. (2024). \textit{PostgreSQL 16 documentation}.
      Recuperado el 10 de noviembre de 2026, de \url{https://www.postgresql.org/docs/16/}

\item Tectonic Typesetting. (s.\ f.). \textit{Tectonic: a LaTeX engine for modern
      workflows}. Recuperado el 10 de noviembre de 2026, de
      \url{https://tectonic-typesetting.github.io/}
\end{list}

%=======================================================================
\newpage
\appendix
\section{Ap\'endice A. C\'odigo completo de la aplicaci\'on}

\begin{lstlisting}[caption=app/app.py,language=Python]
import os
import time
from contextlib import contextmanager

import psycopg2
from flask import Flask, jsonify, request

app = Flask(__name__)

DB_HOST = os.environ.get("DB_HOST", "db")  # "db" = nombre del servicio, NO localhost
DB_PORT = os.environ.get("DB_PORT", "5432")
POSTGRES_DB = os.environ["POSTGRES_DB"]
POSTGRES_USER = os.environ["POSTGRES_USER"]
POSTGRES_PASSWORD = os.environ["POSTGRES_PASSWORD"]


def get_conn():
    """Abre una conexion nueva a PostgreSQL usando el host por nombre de servicio."""
    return psycopg2.connect(
        host=DB_HOST,
        port=DB_PORT,
        dbname=POSTGRES_DB,
        user=POSTGRES_USER,
        password=POSTGRES_PASSWORD,
        connect_timeout=5,
    )


@contextmanager
def conexion():
    """Context manager: commit/rollback automatico y cierre de la conexion."""
    conn = get_conn()
    try:
        with conn:
            with conn.cursor() as cur:
                yield cur
    finally:
        conn.close()


def esperar_base_datos(intentos=30, espera=2):
    """Reintenta la conexion hasta que PostgreSQL acepte consultas."""
    for i in range(intentos):
        try:
            conn = get_conn()
            conn.close()
            print(f"[init] Base de datos disponible en {i + 1} intento(s).", flush=True)
            return True
        except psycopg2.OperationalError as exc:
            print(f"[init] BD no lista ({i + 1}/{intentos}): {exc}", flush=True)
            time.sleep(espera)
    return False


def init_db():
    """Crea la tabla mensajes si aun no existe."""
    with conexion() as cur:
        cur.execute(
            """
            CREATE TABLE IF NOT EXISTS mensajes (
                id     SERIAL PRIMARY KEY,
                texto  TEXT NOT NULL,
                creado TIMESTAMPTZ NOT NULL DEFAULT NOW()
            )
            """
        )


@app.get("/")
def health():
    """Endpoint de salud: comprueba que la app alcanza a la base de datos."""
    estado_db = "desconocido"
    try:
        with conexion() as cur:
            cur.execute("SELECT 1")
            cur.fetchone()
        estado_db = "ok"
    except psycopg2.Error as exc:
        estado_db = f"error: {exc}"

    return jsonify(status="ok", servicio="app", db_host=DB_HOST, db_estado=estado_db)


@app.post("/mensajes")
def crear_mensaje():
    """Guarda un mensaje en PostgreSQL (app -> db)."""
    datos = request.get_json(silent=True) or {}
    texto = (datos.get("texto") or "").strip()
    if not texto:
        return jsonify(error="El campo 'texto' es obligatorio"), 400

    with conexion() as cur:
        cur.execute(
            "INSERT INTO mensajes (texto) VALUES (%s) RETURNING id, texto", (texto,)
        )
        nuevo_id, texto_guardado = cur.fetchone()

    return jsonify(id=nuevo_id, texto=texto_guardado), 201


@app.get("/mensajes")
def listar_mensajes():
    """Devuelve todos los mensajes almacenados en PostgreSQL (db -> app)."""
    with conexion() as cur:
        cur.execute("SELECT id, texto FROM mensajes ORDER BY id")
        filas = cur.fetchall()

    return jsonify([{"id": i, "texto": t} for i, t in filas])


@app.get("/mensajes/<int:mensaje_id>")
def obtener_mensaje(mensaje_id):
    """Devuelve un mensaje puntual por su identificador."""
    with conexion() as cur:
        cur.execute("SELECT id, texto FROM mensajes WHERE id = %s", (mensaje_id,))
        fila = cur.fetchone()

    if fila is None:
        return jsonify(error="Mensaje no encontrado"), 404

    return jsonify({"id": fila[0], "texto": fila[1]})


if __name__ == "__main__":
    if esperar_base_datos():
        init_db()
        print("[init] Tabla 'mensajes' verificada.", flush=True)
    else:
        print("[init] ERROR: la base de datos no respondio; se reintenta al reiniciar.",
              flush=True)

    # 0.0.0.0 es obligatorio para que el contenedor sea alcanzable desde fuera.
    app.run(host="0.0.0.0", port=5000)
\end{lstlisting}

\section{Ap\'endice B. Gu\'ia r\'apida de puesta en marcha}

\begin{lstlisting}[caption=Secuencia de comandos de la Entrega 1,style=estiloshell]
# 1. Crear el .env real a partir de la plantilla
copy .env.example .env        # Windows
cp .env.example .env          # Linux / macOS / WSL

# 2. Construir y levantar
docker compose build
docker compose up -d
docker ps

# 3. Inspeccionar la red creada por Compose
docker network ls
docker network inspect proyecto-ci_red_proyecto

# 4. Probar la comunicacion app <-> db
curl -X POST http://localhost:5000/mensajes \
     -H "Content-Type: application/json" -d "{\"texto\":\"hola\"}"
curl http://localhost:5000/mensajes
docker compose exec app python -c "import socket; print(socket.gethostbyname('db'))"
docker compose exec db psql -U usuario -d proyecto -c "SELECT * FROM mensajes;"

# 5. Limpiar (con -v se borra tambien el volumen)
docker compose down -v
\end{lstlisting}

\section{Ap\'endice C. Lista de capturas requeridas}

\begin{center}
\begin{tabularx}{\textwidth}{l X}
\toprule
\textbf{Archivo en \texttt{figuras/}} & \textbf{Contenido esperado} \\
\midrule
\texttt{01-git-log.png} & Historial de commits del repositorio \\
\texttt{02-docker-build.png} & Salida de \texttt{docker compose build} \\
\texttt{03-docker-compose-up.png} & Salida de \texttt{docker compose up -d} \\
\texttt{04-docker-ps.png} & Salida de \texttt{docker ps} \\
\texttt{05-network-inspect.png} & Salida de \texttt{docker network inspect} \\
\texttt{06-curl-post.png} & \texttt{POST /mensajes} con respuesta \texttt{201} \\
\texttt{07-curl-get.png} & \texttt{GET /mensajes} con el dato insertado \\
\texttt{08-dns-db.png} & Resoluci\'on de \texttt{db} desde el contenedor \texttt{app} \\
\texttt{09-psql-select.png} & \texttt{SELECT * FROM mensajes;} en \texttt{psql} \\
\texttt{10-actions-verde.png} & Ejecuci\'on del pipeline en verde \\
\bottomrule
\end{tabularx}
\end{center}

\end{document}
